% Options for packages loaded elsewhere
\PassOptionsToPackage{unicode}{hyperref}
\PassOptionsToPackage{hyphens}{url}
\documentclass[
  11pt,
]{article}
\usepackage{xcolor}
\usepackage[margin=1in]{geometry}
\usepackage{amsmath,amssymb}
\setcounter{secnumdepth}{-\maxdimen} % remove section numbering
\usepackage{iftex}
\ifPDFTeX
  \usepackage[T1]{fontenc}
  \usepackage[utf8]{inputenc}
  \usepackage{textcomp} % provide euro and other symbols
\else % if luatex or xetex
  \usepackage{unicode-math} % this also loads fontspec
  \defaultfontfeatures{Scale=MatchLowercase}
  \defaultfontfeatures[\rmfamily]{Ligatures=TeX,Scale=1}
\fi
\usepackage{lmodern}
\ifPDFTeX\else
  % xetex/luatex font selection
  \setmainfont[]{Times New Roman}
\fi
% Use upquote if available, for straight quotes in verbatim environments
\IfFileExists{upquote.sty}{\usepackage{upquote}}{}
\IfFileExists{microtype.sty}{% use microtype if available
  \usepackage[]{microtype}
  \UseMicrotypeSet[protrusion]{basicmath} % disable protrusion for tt fonts
}{}
\makeatletter
\@ifundefined{KOMAClassName}{% if non-KOMA class
  \IfFileExists{parskip.sty}{%
    \usepackage{parskip}
  }{% else
    \setlength{\parindent}{0pt}
    \setlength{\parskip}{6pt plus 2pt minus 1pt}}
}{% if KOMA class
  \KOMAoptions{parskip=half}}
\makeatother
\setlength{\emergencystretch}{3em} % prevent overfull lines
\providecommand{\tightlist}{%
  \setlength{\itemsep}{0pt}\setlength{\parskip}{0pt}}
\usepackage{bookmark}
\IfFileExists{xurl.sty}{\usepackage{xurl}}{} % add URL line breaks if available
\urlstyle{same}
\hypersetup{
  hidelinks,
  pdfcreator={LaTeX via pandoc}}

\author{}
\date{}

\begin{document}

\section{Response to Reviewer 1}\label{response-to-reviewer-1}

\textbf{Manuscript:} \emph{Proxy-Based Diagnostics of Quantum Oracle
Sketching Robustness for Non-IID Sensor and Telemetry Streams}
(submitted as \emph{Quantum Oracle Sketching for Non-IID Sensor and
Telemetry Streams: A Proxy-Based Computational Study})\\
\textbf{Journal:} Sensors (MDPI) · \textbf{Manuscript ID:}
sensors-4470240

We thank both reviewers for careful, constructive reviews that have
materially improved the paper. Both reviews converge on the same central
point, and we agree with it: the previous version, although internally
hedged, still presented itself as evidence about quantum performance,
when what it can honestly support is a \textbf{proxy-based diagnostic
and hypothesis-generation framework for correlation-sensitive streaming
analysis}. The revision repositions the manuscript accordingly ---
title, abstract, statistics, figures, and conclusion --- without
removing any experiment or altering any measured result: all stochastic
outputs regenerate bit-identically from the same seeded pipeline (we
verified this explicitly before re-analysis; the verification scripts
ship with the revision).

\textbf{Summary of the principal changes} (full detail under each point
below):

\begin{enumerate}
\def\labelenumi{\arabic{enumi}.}
\tightlist
\item
  \textbf{Retitled} to \emph{Proxy-Based Diagnostics of Quantum Oracle
  Sketching Robustness for Non-IID Sensor and Telemetry Streams}; the
  abstract now opens ``In theory, \ldots{}'', states that no quantum
  algorithm is implemented or simulated, and was shortened to
  \textasciitilde220 words. \emph{(Note on the tracked-changes file: the
  MDPI class places the title and abstract in the document preamble,
  which latexdiff does not mark; both are completely rewritten.)}
\item
  \textbf{New Section 3.3, ``Relationship Between Operational Proxies
  and Quantum Oracle-Sketching Theory''}, separating established theory
  / assumptions / heuristic constructions / scope limitations, with a
  provenance table (Table 1) stating per quantity what is inherited,
  what is introduced, and what is \emph{not} theoretically guaranteed
  --- including the explicit statement that under the framework's own
  r-only accounting no proxy--baseline crossover exists, so the
  τ-penalised landscape is a classical effective-sample stress test, not
  the framework's cost model.
\item
  \textbf{Statistics rebuilt around estimation rather than testing}
  (Experiment 5): Hodges--Lehmann paired-difference estimates with exact
  signed-rank 95\% intervals, matched-pairs rank-biserial correlations
  (interpreted as sign-concordance), p-values demoted to descriptive
  companions, disclosure that no comparison survives Holm correction at
  n = 10 (with the reason), the τ-estimator lattice sensitivity, and an
  explicit statement that the crossover location is a deterministic
  solution whose uncertainty comes from the τ estimator and the
  constants, not from tests.
\item
  \textbf{New sensitivity analyses from existing data} (no new
  stochastic runs): a C--δ grid for the crossover (Table 8: rho* moves
  from 0.38 to beyond the sweep; ordinal reading is the robust one) and
  an encoder/resolution sensitivity table for the real streams (Table:
  10 re-encodings; the NYC-above-Machine-Temperature ordering holds in
  all 10).
\item
  \textbf{New Discussion subsections} on why analytic and empirical
  correlation estimates diverge, and on what the landscape can and
  cannot predict --- including three \textbf{falsifiable hypotheses
  (H1--H3) with explicit refutation criteria}.
\item
  \textbf{New ``Threats to Validity'' section} (construct / internal /
  external / statistical-conclusion) and a promoted, expanded
  \textbf{Future Work} section.
\item
  \textbf{Conclusion rewritten}: primary contribution = the (τ, r)
  landscape; secondary = proxy-based hypothesis generation; the sentence
  ``The paper demonstrates no quantum advantage and implements no
  quantum algorithm'' appears verbatim.
\item
  \textbf{Figures}: the Experiment-5 legend entry ``Proxy-model
  advantage'' renamed ``Proxy-estimated separation''; in-image labels
  ``n = 10 qubits'' / ``Quantum O(n)'' removed from Figures 1, 2, 6;
  captions aligned with the proxy framing; Figure 8(d)'s synthetic
  markers disclosed as schematic; typographic and overfull fixes
  throughout.
\item
  \textbf{Reference audit (preprints and metadata)}: all 54 references
  were classified, and every entry was then verified at the metadata
  level --- all 51 DOIs were resolved through doi.org content
  negotiation and machine-compared against the cited title, first
  author, venue, and year, with the non-DOI entries checked at source.
  Two corrections resulted: (i) Kallaugher, Parekh and Voronova (SOSA
  2025) now carries the SIAM proceedings DOI (10.1137/1.9781611978315.2)
  instead of its arXiv link, with the printed byline order confirmed on
  the publisher's pages; and (ii) a transposed article number in the Su,
  Guo and Wang (2024) entry was corrected (Information Sciences 676,
  article 120799, DOI 10.1016/j.ins.2024.120799 --- the previously cited
  120783 resolves to an unrelated article). One preprint is retained and
  explicitly flagged: Zhao et al.~(arXiv:2604.07639) is the framework
  this paper studies; we verified via the arXiv record, the authors'
  official listing, and a Crossref registry query that no peer-reviewed
  version exists to date, the entry is labelled ``Preprint'', and every
  invocation of its theorems is co-cited with peer-reviewed anchors
  (Kallaugher, FOCS 2021; Huang et al., Science 2022; Kallaugher et al.,
  SOSA 2025). A citation--bibliography consistency check confirms all 54
  in-text keys map one-to-one onto the 54 entries. Following the
  Academic Editor's request that preprints not be cited, that single
  remaining entry now also carries an explicit footnote at its first
  citation in the Introduction, stating its preprint status and why no
  conclusion of this paper depends on its final published form; a
  separate response to the Academic Editor accompanies this letter.
\end{enumerate}

A tracked-changes PDF (\texttt{tracked\_changes.pdf}, latexdiff against
the reviewed version) accompanies this letter.

\begin{center}\rule{0.5\linewidth}{0.5pt}\end{center}

\subsection{Point-by-point response}\label{point-by-point-response}

We thank the reviewer for the structured list of fifteen
recommendations; every one is addressed below.

\textbf{R1.1 --- ``I strongly recommend to reduce the emphasis on
quantum advantage and instead presenting the τ,r landscape primarily as
a diagnostic framework to identify scenarios where future quantum
implementations is most promising.''}

\emph{Response.} Adopted as the organising principle of the revision.
The title now leads with ``Proxy-Based Diagnostics''; the Introduction
states the paper ``is best read as a diagnostic and
hypothesis-generation framework \ldots{} it neither demonstrates nor
claims a realised quantum--classical separation'' (§1); the Conclusion
names the landscape as the primary contribution and hypothesis
generation as the secondary one; and §7.4 now states three falsifiable
hypotheses (H1--H3) that make ``where future implementations are most
promising'' concrete and testable. \emph{Changes:} Title; Abstract; §1
(scope paragraph, contributions 3--4); §7.4 (Scope of the Diagnostic,
incl.~H1--H3); §10 (Conclusions, rewritten); every
``quantum-favourable'' replaced by ``proxy-favourable''/``hypothesized
favourable''.

\textbf{R1.2 --- ``Please clarify the relationship between the heuristic
proxy model and the original theoretical framework.''}

\emph{Response.} A dedicated subsection now does exactly this. §3.3
separates (i) \emph{established theory} --- the ×R sample overhead,
O(n)-qubit memory, and task-specific O(N\^{}c) classical hardness, all
theorems of the original framework; (ii) \emph{assumptions} this paper
introduces; (iii) \emph{heuristic constructions} --- the additional τ
division in T\_eff (deliberately more pessimistic than the framework's
own accounting, under which no crossover would exist at all) and the
VC-type task rate used in place of oracle-construction cost; and (iv)
\emph{scope limitations}. Table 1 tabulates the provenance and epistemic
status of every model-level quantity, with a ``not guaranteed'' column.
\emph{Changes:} New §3.3 with four labelled paragraphs and Table 1; §3.2
proxy-status paragraph.

\textbf{R1.3 --- ``The analytic Markov τ proxy differs substantially
from the empirical measurements. I recommend to discuss this discrepancy
in greater detail and explain whether it reflects limitations of the
proxy model or the synthetic data generation process.''}

\emph{Response.} A dedicated subsection (§7.3, ``Why Analytic and
Empirical Correlation Estimates Diverge'') now answers the question
directly. For Markov switching the two quantities answer different
questions: each \emph{bit's} marginal autocorrelation decays as ρ\^{}k
and crosses 1/e at ≈2.8 lags (ρ = 0.7), while the \emph{joint} n-bit
chain takes Θ(n log n/(1−ρ)) steps to mix --- both numbers are correct
answers to different questions. The divergence is therefore partly
expected behaviour and partly a limitation of single-scalar proxies; it
is \textbf{not} a data-generation defect: the generators reproduce their
intended statistics, and the gap persists under all four alternative τ
estimators of Experiment 9, isolating the estimator family as the cause.
The subsection also derives the practical implications (empirical τ as
the operational input; analytic τ as a conservative envelope; their
ratio as a long-memory diagnostic). \emph{Changes:} New §7.3 (three
labelled paragraphs: technical explanation / interpretation /
implications); Experiment 4 cross-reference; Conclusion RQ1.

\textbf{R1.4 --- ``I recommend to validate the proposed framework on
additional real-world datasets. It will help to better demonstrate its
generalizability.''}

\emph{Response.} We agree this is the most valuable extension, and we
are explicit that we have not done it within this revision: adding
datasets would have been a new experimental campaign, and we preferred
to keep this revision strictly re-analysis-based so that every published
number remains bit-reproducible. What we did add, from the same raw
data, is an \textbf{encoder/resolution sensitivity analysis} (new table
in §6.8): the two streams were re-encoded across ten sampling-resolution
and bucket-width configurations, and the
NYC-Taxi-above-Machine-Temperature proxy ordering holds in all ten ---
so the existing two-stream placement is at least not an artefact of one
encoder setting. The manuscript now states plainly that the real-stream
experiment ``is illustrative rather than statistically representative''
(§6.8), the Threats section records the n = 2 limitation, and Future
Work item (iii) commits to network telemetry, cybersecurity event
streams, financial series, and IoT fleets, with hypothesis H1's protocol
(rank correlation between placement and outcome over a ≥6-stream corpus)
defining exactly how generalisability will be tested. \emph{Changes:}
§6.8 (illustrative statement; encoder/resolution table and paragraph);
§8 External validity; §9 Future Work (iii); §7.4 H1--H2.

\textbf{R1.5 --- ``Please include a more comprehensive sensitivity
analysis of the selected proxy parameters, particularly C and δ, by this
you will show how they influence the reported results.''}

\emph{Response.} Done, quantitatively. Because the proxy is a
closed-form function of the measured (τ, r), we re-evaluated it on the
same per-seed empirical τ values across C in \{1, 1.5, 2, 2.5\} crossed
with δ in \{0.01, 0.05, 0.1\}. New Table 8 reports the induced movement
of the proxy--baseline crossover: regime \emph{ordering} is invariant (C
and δ enter only through the monotone factor κ = C²log(1/δ)), but the
crossover location moves from ρ* ≈ 0.38 at (C = 2.5, δ = 0.01) to beyond
the sweep range at C = 1, with ρ* ≈ 0.82 at the operating point. The
text now instructs that the landscape be read \textbf{ordinally}, the
abstract carries the conditionality explicitly, and we additionally
state that C = 2 is a conventional order-unity choice not fixed by any
constant in the cited theory (§5.3). \emph{Changes:} New ``Proxy
constants C and δ'' paragraph + Table 8 in §6.9; §3.2 and §5.3 updated;
Abstract; §8 Statistical-conclusion validity.

\textbf{R1.6 --- ``The comparison between deterministic quantum proxies
and stochastic classical baselines should be presented more cautiously.
The current statistical analysis may unintentionally suggest a
comparison between implemented algorithms rather than between empirical
results and theoretical proxies.''}

\emph{Response.} We rebuilt the statistical presentation around this
point. Table 4 now reports Hodges--Lehmann estimates of the paired
(proxy − classical) difference with \textbf{exact} signed-rank 95\%
intervals and rank-biserial correlations; p-values are retained only as
descriptive companions. A dedicated paragraph (``What these tests do and
do not establish'', §6.5) states that the pairing is
stochastic-learner-versus-model-prediction, that the proxy's seed
variability comes solely from the per-seed τ estimate, that r\_rb here
indexes \emph{sign concordance across seeds} rather than magnitude, and
--- verbatim --- that the tests ``do not, and cannot, establish that one
algorithm outperforms another, and no quantum algorithm is implemented
anywhere in this comparison''. We further disclose that none of the five
comparisons survives Holm correction at n = 10 (smallest adjusted p =
0.098; an outcome of one-to-two discordant seeds, not a design ceiling),
that the ρ = 0.48 interval is not robust to a single step of the
integer-valued τ estimator, and that the crossover location is the
deterministic solution of acc\_proxy(τ̂(ρ)) = mean classical accuracy,
with uncertainty inherited from τ̂ and (C, δ), not from tests.
\emph{Changes:} Table 4 (rebuilt, incl.~caption stating the five
pre-specified test points); §6.5 body and interpretation paragraph;
Abstract (p-value removed); §7.1; §8 Statistical-conclusion validity.

\textbf{R1.7 --- ``I recommend to provide a stronger justification for
the selected classical baselines and discussing if additional streaming
methods can offer a more informative comparison.''}

\emph{Response.} §5.2 now justifies the trio explicitly: the three
baselines span complementary axes of bounded-memory streaming learning
--- gradient-based (online SGD), variance-reduced gradient-based
(averaged SGD), and non-gradient hash-count-based (Count-Min) --- while
sharing the single-pass, fixed-memory regime that the proxy models. We
also explain why drift-adaptive methods (ADWIN-equipped learners,
adaptive random forests) are deliberately excluded from the primary
comparison --- they respond to non-stationarity by discarding history,
which would confound the correlation-budget question the proxy isolates
--- and we name benchmarking them on the same landscape as an extension.
\emph{Changes:} §5.2 (new justification passage with citations); §9
Future Work.

\textbf{R1.8 --- ``Please explain more clearly why the proposed
operational proxies, τ and r, are sufficient to characterize the
correlation properties of the studied data streams, and discuss any
limitations of relying only on these two measures.''}

\emph{Response.} The revision treats two-scalar sufficiency as an
explicit \emph{assumption}, not a fact: §3.3 (``Assumptions'', item (a))
states it as a modelling choice of this paper, and §7.3 shows its
precise failure mode --- for non-mixing (long-range-dependent) processes
no single scalar exists that the analytic and empirical estimators could
agree on, so a single-number τ is intrinsically lossy there. The
mitigations are stated: the analytic-to-empirical τ \emph{ratio} is
itself a useful third diagnostic flagging long-memory content, and
multi-scale refinements of τ (integrated-autocorrelation estimates at
several horizons) are named as follow-on work. The Threats section
(Construct validity) records the residual risk. \emph{Changes:} §3.3
Assumptions; §7.3 Interpretation and Implications; §8 Construct
validity; §9 Future Work (ii).

\textbf{R1.9 --- ``I recommend to expand the discussion of the
assumptions behind each synthetic correlation regime and to please
explain how closely these settings represent practical sensor and
telemetry environments.''}

\emph{Response.} §4 already carried per-regime provenance notes; the
Summary (§4.6) now adds the realism mapping: Markov switching idealises
short-range state persistence in sensor readouts; seasonal drift the
diurnal/weekly cycles of telemetry and demand series; burst repetition
the duplicate-heavy event bursts of monitoring and logging pipelines;
and long-range dependence the self-similar traffic behaviour first
measured on Ethernet --- with the explicit caveat that these are
idealisations, that real streams mix the mechanisms (as the two NAB
cases show), and that the mapping to any concrete deployment is
qualitative (Table 9). \emph{Changes:} §4.6 (new passage with
citations); Table 9 caption (analytic-convention statement).

\textbf{R1.10 --- ``The manuscript would benefit from a dedicated
discussion of threats to validity, include the assumptions introduced by
the heuristic proxy model, synthetic data generation and feature
binarization.''}

\emph{Response.} Added as a dedicated section. §8 (``Threats to
Validity'') is organised by construct, internal, external, and
statistical-conclusion validity and covers exactly the requested items:
the heuristic proxy constructs (with pointers to the provenance table),
synthetic-generator selection, feature binarisation (arcsine shrinkage
makes real-stream (τ, r) estimates lower bounds on the
continuous-process correlation), the deterministic-vs-stochastic
comparison design, dataset selection, target definition (including the
full-series threshold look-ahead raised in review), fixed
hyperparameters, and the n = 10 seed budget. \emph{Changes:} New §8
replacing the former limitations list.

\textbf{R1.11 --- ``I recommend to discuss the computational complexity
of estimating τ and r and the scalability of the proposed framework for
high-volume streaming systems.''}

\emph{Response.} Added. §5.1 now contains an ``Estimator cost and
scalability'' paragraph: the τ estimator needs the bitwise
autocorrelation up to lag K --- O(nTK) time naively, O(nK) memory with a
ring buffer per feature --- and the r estimator scans a W-length forward
window --- O(TW) comparisons, expected O(T) with hashing, O(Wn) memory.
Both are linear in T, independent of the alphabet size N = 2\^{}n,
single-pass friendly, and embarrassingly parallel across features; on
the longest stream studied (T = 22,695) they complete in well under a
second, so the diagnostic adds negligible overhead for high-volume
deployments. \emph{Changes:} §5.1 (new paragraph).

\textbf{R1.12 --- ``Some figures contain many overlapping curves and
annotations, I recommend to simplify the visual presentation and enlarge
labels to improve readability.''}

\emph{Response.} All nine figures were re-audited. Concretely: the
Figure 5 legend entry ``Proxy-model advantage'' was renamed
``Proxy-estimated separation'' and stray typographic backslashes were
removed from five panel titles (Figures 4--6 regenerated); in-image
labels ``n = 10 qubits'', ``Quantum proxy (IID, 10 qubits)'', and
``Quantum O(n)'' were replaced with neutral proxy/cited-theory labels
(Figures 1, 2, 6 regenerated); Figure 8(d)'s synthetic markers are now
labelled as schematic orientation marks and the caption gives the
numeric stream coordinates instead of a visual-adjacency claim; caption
CI statements were scoped to the stochastic panels only; and four tables
that previously overran the text width were reformatted. All
regenerations are rendering-only --- the numerical outputs are
byte-identical. (The landscape figures already use colour-blind-safe
palettes with direct contour labelling from the previous revision
cycle.) \emph{Changes:} Figures 1, 2, 4, 5, 6 regenerated; captions of
Figures 1, 3, 5, 6, 7, 8, 9; Tables 2--9 formatting.

\textbf{R1.13 --- ``The discussion could better distinguish theoretical
quantum memory advantages from the practical observations obtained
through the proxy-based evaluation.''}

\emph{Response.} This distinction is now enforced at every occurrence.
Experiment 6 opens: ``No memory lower bound is measured empirically in
this experiment'', and states that no result constitutes an empirical
memory floor for online SGD, averaged SGD, or Count-Min; §3.3 (``Memory
references'') makes the same statement structurally; the provenance
table lists the Ω(√N) curve as a ``conservative theoretical reference''
whose empirical-floor status is explicitly not guaranteed; Figure 6's
caption is labelled ``cited bounds; no measured content''; and §7.1
closes its three structural factors with ``All three factors are
properties of the proxy model and the cited theory; none is a measured
quantum result'', with the n ≥ 20 statement marked as an extrapolation
beyond the tested range. \emph{Changes:} §3.3; §6.6 (opening and panel
text); Figure 6 caption; §7.1; Table 1 row.

\textbf{R1.14 --- ``Please discuss the applicability of the proposed
framework to other streaming domains beyond sensor telemetry, such as
network traffic, financial streams, or cybersecurity monitoring.''}

\emph{Response.} Table 9 maps network backbone traffic, financial
returns, sensor/IoT, and concept-drifted streams onto the landscape from
published measurement ranges (now recomputed from the paper's own
formula, with verdicts given as ranges and the estimator convention
stated), and Future Work item (iii) commits to exactly the requested
domains: network telemetry (flow and latency series), cybersecurity
event streams (intrusion and log-anomaly feeds), financial tick and
returns series, and large-scale IoT monitoring fleets. \emph{Changes:}
Table 9 (recomputed labels and caption); §7.2; §9 Future Work (iii).

\textbf{R1.15 --- ``strengthen the conclusion, for this emphasize the
practical contributions of the study, clearly summarize its limitations,
and please outline the most important directions for future research.''}

\emph{Response.} The Conclusion is rewritten to exactly this structure:
a ``Practical contribution'' paragraph (the landscape as a cheap
pre-screening step --- estimate (τ, r), place the stream, read off
whether correlation alone erodes the budget below the
rare-event-learning level, as it does for Machine Temperature at T\_eff
= 52); a ``Limitations and required validation'' paragraph pointing to
the Threats section and naming the decisive missing step (an implemented
or simulated oracle-sketching pipeline evaluated at predicted-favourable
and predicted-unfavourable coordinates); and research-question answers
restated in effect-size terms with estimator conventions bound.
\emph{Changes:} §10 (full rewrite); §9 Future Work.

\begin{center}\rule{0.5\linewidth}{0.5pt}\end{center}

\subsection{Note on verification}\label{note-on-verification}

Before making any statistical change we reproduced the published
Experiment-5 per-seed results bit-exactly from the pinned seeds (all
five Wilcoxon p-values match the published values to six significant
digits), and all figure regenerations are rendering-only (numerical
outputs byte-identical). The re-analysis and sensitivity scripts, with
their outputs, are included in the accompanying package
(\texttt{verification\_scripts/exp5\_reanalysis.py} and
\texttt{verification\_scripts/encoder\_sensitivity.py};
\texttt{code/exp5\_effectsizes.py} in the public repository), and the
reproducibility documentation was corrected and extended (exact seed
indices per experiment; the previously missing Machine-Temperature
download step).

On behalf of all authors,

\textbf{AbdelMoniem Helmy} (corresponding author)
abdelmoniem.hafez@cu.edu.eg · ORCID 0000-0001-5996-6019

\end{document}
